\cotefile{156-4}
\scanpage{18}\mspage{17}
% @ opening
Une \uline{\unsure{contrée}} est définie par

% @ 1°
1°) Un ensemble de ``figures élémentaires'', $\unsure{F_{el}|F_{\mathrm{el}}}$

% @ 2°
2°) Une relation d'ordre $\ll$ \deleted{$\unsure{\beta|b}$} (raffinement) parmi les figures élémentaires

% @ 3°
3°) Une relation \unsure{réflexive} symétrique $\unsure{R|\mathcal{R}}$ (\deleted{\unsure{admissibles|admissibilité|constructibles}}\inserted{``compatibilité''}) parmi figures élémentaires

% @ NB, oblique from left across 1°–4°
% slanted
\margin{\uline{NB} Il revient au $=$ de se donner 4°), ou de se donner la notion de \uline{figure} [une \deleted{\unsure{partie}} \unsure{configuration} est \unsure{une} figure \unsure{si}, en \unsure{même} \unsure{tps} que \unsure{chaque|toute} figure él. qu'elle \unsure{contient}, elle \unsure{contient} \unsure{aussi} \unsure{toute} \unsure{figure} $\leq$ \unsure{une} qu'elle \unsure{englobe|engendre|contient}] et la \unsure{propriété} pour une figure \unsure{d'être} \unsure{finie}.}

% @ 4°
4°) Un ens de \unsure{parties|partie} de $F_{el}$, appelées les \unsure{---} finies -- dans le cas d'une ``\unsure{contrée} \uuline{finie}'' ce sont les parties \uline{finies} qui sont des ``configurations'' \unsure{ou} i.e. \unsure{familles|famés} de figures él. deux à deux compatibles
% @ small note at left of 4°
% slanted
\margin{\uline{contrée} \uline{finie} ou plutôt, ``\uline{compacte}'' : celle où toute figure est finie}
Satisfaisant à des axiomes convenables, à dégager. Il y a déjà celui-ci

% @ condition 1
\uline{\unsure{Contr|Cond}}\textsubscript{1} \deleted{Si} Si \deleted{\unsure{p}}\unsure{parmi|pour} deux fig. élémentaires $\deleted{\unsure{Z}}X, X'$, on \unsure{pose}
\[
\begin{aligned}
X'\leq X\quad &\text{ssi (déf)}\quad X'\ll X,\ \text{et }X'\text{ et }X\\
&\hphantom{\text{ssi (déf)}\quad X'\ll X,\ }\text{sont comp.}
\end{aligned}
\]
alors on trouve une relation \uline{d'ordre},

% @ Figs 1 and Fig 2
De plus, en termes de cette relation d'ordre, on a les propriétés \uline{\unsure{Figs|Fig}} 1

\margin{$\vert\vert\vert$}
$(\Longleftrightarrow)$ [Si $X', X''\ll X$, \inserted{et si $\unsure{X|X'},\unsure{X'|X''}$} sont compatibles à $X$, ils sont compatibles entre eux])

\uline{Fig}\textsubscript{2} (dim combinatoire finie) On va
\scanpage{19}\mspage{18}
% @ first line
la \deleted{\unsure{raffiner en}} \unsure{renforcer}. \deleted{\illegible{1}}

% @ first, crossed condition
\deleted{\uline{\unsure{Contr|Cond}}\textsubscript{2} Si $X\in\unsure{F_{\grave e l}|F_{\acute e l}|F_{el}}$, alors l'ens \unsure{des} \unsure{$X'\ll X$}}
\[
\deleted{\{X'\in F_{el}\mid X'\ll \unsure{F_{el}|X}\}}
\]

% @ second condition
\uline{Contr}\textsubscript{2} Pour \unsure{tt} $X\in F_{\bar e l}$, considérons la dim comb. (pour $\leq$) des $X'\in F_{el}$ \deleted{tels que} qui raffinent $X$. Alors les dim. comb. de ces $X'$ sont majorées.

% @ dimension, middle
\unsure{Leur} borne sup. s'appelle la \uline{dimension \inserted{(\unsure{géom})} de $X$}. \deleted{\unsure{U}}La dim géom d'une figure, est par déf. le sup des dim géom. des \deleted{\unsure{ses}} figures élémentaires qui \unsure{la|les} constituent.
% slanted
\margin{\deleted{\unsure{Donc|D'où} la dimension}}

% @ interline chain
\inserted{C'est donc le sup des longueurs des chaînes $(X_0,X_1,\ldots,X_n)$ avec
\[
 X_0<X_1\ \ldots<X_{n-1}\ll X_n\in F.
\]}
% @ right edge below chain
\margin{\unsure{Elle|Ell} \deleted{$\unsure{F}$} peut être infinie pour $F$ \unsure{infinie}. Si $F'$ raffine $F$, $\unsure{\dim F'|\dim F}$ $\unsure{\leq|=}$ $\dim F$}

% @ dated paragraph
\margin{(\uline{11 juin})}
On aurait pu développer \unsure{un peu|un jour|une fois} les notions de ``partie constructible'', et adhérence, intérieur, frontière, composantes connexes associés à une figure, \unsure{avec} les \unsure{seules} \unsure{données|donnés} $\leq$, compatibilité sur \deleted{\unsure{la}} \deleted{\unsure{un}} $F_{\text{él.}}$.
% @ lower oblique note
% slanted
\margin{Aussi, la notion de \uline{structure de \unsure{contrée}} induite sur une \uline{\unsure{multistrate|multistr.}} figure. La notion de \unsure{sous-figure|sous figure} ($\unsure{\leq}$), \uline{composantes} \unsure{d'un} \unsure{univers|ensemble} \unsure{des} \unsure{points|parties} \unsure{pour} \unsure{``disj.''} \unsure{d'une} \unsure{figure}, \unsure{connexes|maximales}, celle \unsure{de|des} \illegible{1}, \unsure{Aussi}.}
% @ closing ordinary lines
J'y reviendrai sûrement. Pour le moment, j'ai envie de pousser \deleted{\unsure{plutôt}} un peu \unsure{ds|dans} la direction $\ll$, compatibilité,
\scanpage{20}\mspage{\unsure{19|17}}
% @ first three lines
pour arriver à la notion de lieu, et d'interprétation des données en termes de lieux.

% @ terminology
\unsure{Pt|P\textsuperscript{t}} de vue \unsure{terminologie|terminologique}, au lieu de ``figure élémentaire'', je préfère dire maintenant ``\uline{multistrate}'', plus parlant.

Donc une figure quelconque est un ens de multistrates mutuellement compatibles, fermé pour $\leq$, et loc. fini.

% @ definition of lieu
\deleted{\unsure{F|E}} On appelle \deleted{\unsure{multi}} \uline{lieu} de la contrée, une multistrate qui est minimale pour $\ll$, i.e. qui n'a d'autre raffinement qu'elle-\unsure{\={m}|m}. \deleted{C'est donc} \deleted{\unsure{point|\c{c}a} \unsure{d'un pt}} \deleted{C'est donc} \deleted{Dans le} (Dans les exemples-types, c'est un ens de parties de $\deleted{\unsure{A}}X$, réduits à une seule partie, qui est \unsure{ponctuelle} -- \unsure{donc|dans} les lieux sont les p\textsuperscript{ts} de $X$ \unsure{..|.})

% @ multistrates minimales
Ce sont donc aussi des multistrates minimales pour $\leq$, \deleted{\unsure{donc}} \inserted{à fortiori} des ``\unsure{variétés|vacuités}'' \deleted{\unsure{connexes} \unsure{non} \unsure{vides})''} -- ce sont les non-variétés \unsure{0-connexes|o-connexes} qui n'ont pas de raffinement autre qu'elle-\={m}.

% @ last lines and inclusion
Soit $\unsure{L|\mathcal{L}}$ l'ensemble des lieux, \deleted{Pour tout} donc
\[
\begin{aligned}
L\subset M\quad (={}&F_{el}=\text{ensemble des}\\
&\hphantom{F_{el}={}}\text{multistrates}\,)
\end{aligned}
\]
\scanpage{21}\mspage{20}
% @ upper lines, L_X
\deleted{toute} Pour toute multistrate $X$, soit $\unsure{L|\mathcal{L}}_X$ \inserted{$\in\unsure{\mathcal P|P}(L)$} l'ens des lieux qui raffinent $X$.
\margin{$\widehat L_X$}
et soit $\widehat L_X\subset\mathcal P(X)$ i.e $\widehat L_X\in\mathcal P(\mathcal P(X))$ l'ensemble des \deleted{$\unsure{X|Y}$} $L_Y$ pour $Y\leq X$ ($Y$ sous-multistrate de $X$). \deleted{Considérons} \deleted{les}

% @ structures and map
J'aimerais réconstituer les structures $\ll$ et comp. sur $M$, en termes de l'application
\[
 X\longmapsto\widehat L_X\qquad M\longrightarrow\mathcal P(\mathcal P(X)).
\]

% @ figs, label 3
\deleted{\uline{\unsure{Contr|Cond}}\textsubscript{3}} \unsure{Itou|Alors} pour $\unsure{\mathcal F|F}\unsure{,|.|{}_{e}}\subset\mathcal P(M)$, (ens des figures), \unsure{en|on} \unsure{posant|poserait}
\[
\begin{aligned}
\deleted{\unsure{\widehat F\mapsto|F\mapsto}}\quad &{\deleted{\unsure{\wedge}}L}_F=\text{ens des lieux qui raffinent }F\\
 &\text{i.e.}\quad L_F=\bigcup_{X\in F}L_X\unsure{\subset|\in} L\\
 &\widehat L_F=\{L_X\mid X\in F\}\in\mathcal P(L)\\
 &\hspace{9em}\text{i.e.}\ \in\mathcal P(\mathcal P(L))
\end{aligned}
\]
% @ mapping for figures
qui nous donne
\[
 F\longmapsto\widehat L_F\qquad \mathcal F\longrightarrow\mathcal P(\mathcal P(X))
\]

% @ cancelled condition, foot
\deleted{\unsure{Soient|tout}} \deleted{\uline{Contr}\textsubscript{3} \unsure{a)|i)} Pour que $X\ll Y$ \unsure{soient} $X,Y\in$}

\deleted{Soient $X,Y$ deux multistrates. \unsure{Si} a) Supposons $X\ll Y$. Alors pour tout $X_i\leq X$, \unsure{consi|aussi|on}}
\scanpage{22}\mspage{21}
% @ first application
\deleted{Soit $X$ une multistrate.} L'application
\[
 X_{\unsure{i|{}}}\longmapsto\unsure{L|\mathcal{L}}_{X_i}\qquad
 M_{\leq X}\longrightarrow\unsure{\mathcal P|P}(L)
\]
est bien \unsure{sur|sûr} croissante ($X_{\unsure{j|i}}\leq \unsure{X_i|Y_i}\Longrightarrow X_j\ll \unsure{X_i|Y}\text{ donc }L_{X_j}\subset L_{X_i}$) On aimerait qu'elle le soit \uline{strictement}, i.e. que \deleted{$\unsure{X_j|X'}$}
% @ strict growth
\[
\begin{aligned}
 X_j<X_{\unsure{i|j}}&\Longrightarrow L_{X_j}\subsetneq L_{X_i}\qquad\text{i.e. }\exists\ \text{lieu}\\
 &\qquad \unsure{x|\alpha}\ \text{tel que }x\ll X_i\\
 &\hspace{8em}\text{et }x\not\ll X_j
\end{aligned}
\]

% @ definition of open-circle L
\deleted{\uline{\unsure{NB}}} Soit $\unsure{L_X^{\circ}|L_{X_i}^{\circ}}=L_X\setminus\bigcup_{X_i<X}L_{X_i}$. Condition plus forte
\margin{\uline{\unsure{Contr|Cont}} 3}
\[
\begin{aligned}
 &\forall X\in M_{\deleted{\unsure{X_j|X}}},\quad L_X^{\circ}\ne\varnothing\ \unsure{\text{--}|.}\\
 \text{i.e.}\qquad &\forall\deleted{\unsure{\varnothing|X}}X\in M,\ \exists x\in L,\ \inserted{\text{tel que}}\ x\ll X,\ \text{et}\\
 &\hspace{11em}x\,\unsure{\overline{\leq}|\nleq|\leq}\,X_i\ \text{pour tts}\\
 &\hspace{14em}X_i<X
\end{aligned}
\]

% @ large cancelled block
\deleted{Ainsi, si $X\in M$, l'application}
\[
 \deleted{X_i\longmapsto L_{X_i}^{\circ}\qquad M_{\leq X}\longrightarrow\mathcal P(L)}
\]
\deleted{serait Soient $X$ et $Y$. Soient $X$ et $Y$ deux multistrates. Supposons \unsure{$X\ll$} les compatibles. Je dis}

% @ final paragraph
Soit $X\in M$, considérons les $L_{X_i}^{\circ}$, où $X_i\leq X$. Je dis que leur réunion \inserted{(qui est $\subset L_X$)} est \uline{égale} à $L_X$. En effet, soit $x\in L_X$, donc $\deleted{\unsure{x}}\exists\unsure{i|0}$, $x\in L_{X_i}$, \deleted{et les $\unsure{X_i|Y_i}$ forment un ouvert de $M_{\leq X}$} Il y a un $X_i$ minimal parmi les
\scanpage{23}\mspage{22}
% @ continuation from scan 22
$X_i$ tels que $x\in\unsure{L|\mathcal{L}}_{X_i}$ (à cause \unsure{dim|dem} \unsure{comb|comb.} de $X$ finie) donc $x\in L_{X_i}^{\circ}$. De façon précise, on a
\[
\begin{aligned}
 x\in L_{X_i}^{\circ}&\Longleftrightarrow X_i\text{ est minimal}\\
 &\hspace{3em}\text{parmi les }Y\in M_{\leq X}\\
 &\hspace{3em}\text{tels que }x\in L_Y
\end{aligned}
\]

% @ unique and inequality
J'aimerais que cet $X_i$ soit \uline{unique} i.e. qu'il soit un plus petit élément parmi les $Y$ envisagés donc
\[
 X_i,X_j\deleted{\unsure{X}}\leq X,\quad X_i\ne X_j\Longrightarrow L_{X_i}^{\circ}\cap L_{X_j}^{\circ}=\varnothing
\]
\deleted{\unsure{D}}Donc que $L_X$ soit \unsure{réunion|réunies} \uline{\unsure{disjointe|disjoints}} des $L_{X_i}^{\circ}$ pour $X_i\leq X$. Plus généralement, on aimerait que l'analogue soit vrai pour tte figure $F\subset M$, \deleted{donc} que $L_{\deleted{\unsure{X}}F}$ soit réunion \uline{disjointe} des \deleted{$\unsure{L_X|L_{X_i}}\ \unsure{L_F}$} $L_{X_i}^{\circ}$ pour $X_i\in F$ (alors qu'on sait \inserted{déjà} que c'est la réunion des $L_{X_i}^{\circ}$) ou encore que $\forall x\in L_F$, $\exists$ pl\uline{us petit} $X\in F$, \unsure{telle|tels} que $x\in L_X$ \deleted{\unsure{\={e}lém. min} $X$ tel que $\unsure{x\in L_X}$},

% @ condition numbered 4
\uline{\unsure{Conf|Contr}} 4. Si $X,Y$ sont deux multistrates \uline{compatibles} (donc \deleted{définissent une} \inserted{\unsure{incluses|inclues|entrent} \unsure{dans|ds}} une figure) et $X\ne Y$, alors $L_X^{\circ}\cap L_Y^{\circ}=\varnothing$ (en d'autres termes, si $x\in L_X\cap L_Y$, alors \deleted{on a} $\exists\unsure{X'|Z'}<X$ tel que $x\in L_{X'}$ et \deleted{$\unsure{\exists|X}$} $Y'<Y$ tel que $x\in L_{Y'}$,) \deleted{Donc}
% @ oblique formula at left
% slanted
\margin{\unsure{équivaut|équivmt} : $L_X\cap L_Y=\displaystyle\bigcup_{\substack{Z\ \text{tels}\\\text{que }Z\leq X\\Z\leq Y}}L_Z$ (si $X,Y$ compatibles)}

% @ crossed lower condition
\deleted{\uline{\unsure{Conf|Crt}} 5 Mais \unsure{on|ce} \unsure{veut|que} \inserted{\unsure{aussi|suggère}} \unsure{la réciproque|réciproquement|l'expérience} : si $X,Y\in M$, \unsure{sont} tels que \deleted{\unsure{toutes|toute}} $\forall X'\leq X,\ Y'\leq Y$, avec $X'\ne Y'$, on ait $L_{X'}^{\circ}\cap L_{Y'}^{\circ}=\varnothing$, alors $X'$ et $Y'$ sont compatibles}
\scanpage{24}\mspage{23}
% @ opening and heading
\deleted{\unsure{S|§}} En utilisant \uline{\unsure{Conf|Contr}} 3, \uline{\unsure{Cont|Contr}} 4, on trouve le

\uline{Scholie} Soit $F$ \inserted{$\subset M$} une figure. \deleted{L'application}
\[
\deleted{X\longmapsto\unsure{L|\mathcal{L}}_{\unsure{X|{}}}}
\]

% @ a), parts and partition
a) Les \deleted{ens \unsure{des}} \inserted{parties de $L$} (``strates ouvertes de $F$'') $L_X^{\circ}$ ($X\in F$) sont non vides et mutuellement disjointes, elles forment \deleted{\unsure{de}} une partition de $L_{\deleted{\unsure{X}}F}=\bigcup_{X\in F}L_X=\bigcup_{X\in F}L_X^{\circ}$.

\deleted{b) L'applic} A fortiori, l'application
\[
 X\longmapsto L_X^{\circ}\qquad F\longrightarrow\unsure{\mathcal P|P}(L)
\]
(\unsure{injective|injection})
% @ oblique caveat
% slanted
\margin{Mais bien sûr, si on ne se \unsure{restreint|restreignait} \unsure{pas} aux multistrates qui font partie d'une figure donnée, la connaissance des $L_X^{\circ}$ ou $L_X$, ni même celle des $(L_X,L_X^{\circ})$, ne détermine $X$ (avec sa \uline{multi}-structure, \unsure{i.e.|ou} \unsure{ses} \unsure{parties} \unsure{fermées|ouvertes})}

% @ b), two maps
b) L'application
\[
 X\longmapsto L_X\qquad F\longrightarrow\mathcal P(L)
\]
est également injective. \inserted{[donc si $X,Y$ sont comp. et $L_X=L_Y$ on a $X=Y$)} On a de plus, pour $X,Y\in F$, équivalence entre : les conditions suivantes
% @ four conditions, right margin
\[
\begin{aligned}
 \text{(i)}\quad &X\leq Y\\
 \text{(ii)}\quad &L_X\subset L_Y\\
 \text{(iii)}\quad &L_X^{\circ}\subset L_Y\\
 \text{(iv)}\quad &L_X^{\circ}\cap L_Y\ne\varnothing
\end{aligned}
\]
\margin{donc ces équivalences valent dès que $X,Y$ sont des multistrates \uline{compatibles}.}

% @ c) and right margin
c) Si $X,Y\in\deleted{\unsure{X}}F$, on a
\[
 L_X\cap L_Y=\bigcup_{\substack{Z\in F\\ Z\leq X,\,Z\leq Y}}L_Z
\]
\margin{\unsure{Itou|Idem} pour c)}

% @ proof at foot
Dém de b). \unsure{Les} implications (i)$\Rightarrow$(ii)$\Rightarrow$(iii)$\Rightarrow$(iv) triviales, prouvons (iv)$\Rightarrow$(i) Comme $L_Y=\bigcup_{\unsure{Y_i|Z_i}\leq Y} L_{Y_i}^{\circ}$, \deleted{\unsure{on a}} $\exists Y_i$ avec $L_X^{\circ}\cap L_{Y_i}^{\circ}\ne\varnothing$ donc par a) on a $X=Y_i$ donc $X\leq Y$, \unsure{qcd|qed}
\scanpage{25}
\mspage{24}
% @ opening
Mais, \deleted{\unsure{on}} on arrive aux réciproques

% @ upper left, oblique
% slanted
\margin{Attention, Ces réciproques valent dans les exemples-types \unsure{unistrates|uni-strates} (théorie topologique, ou modérée) mais pas pour les réseaux p.ex.}

% @ Cont 5
\uline{\unsure{Cont|Conf}} 5 Soient $X,Y\in M$ deux multistrates, et considérons $\deleted{\unsure{L_X|\widehat L_X}}\widehat{\unsure{L|\mathcal{L}}}_X,\widehat L_Y\subset\unsure{\mathcal P|P}(L)$ les ens de parties de $L$ correspondantes. Ils permettent, grâce aux scholies de récupérer les strates ``ouvertes''

% @ a), first condition
a) Pour que $X\ll Y$, il f. et s. que \deleted{\unsure{tout} $\forall A\ \unsure{L}\ \forall A\in L_X$} toute strate ouverte de $X$ soit contenue dans une strate ouverte de $Y$

% @ left, beside (a) and (b)
% slanted
\margin{Il en résulte l'énoncé correspondant, pour deux figures $F,F'$, \unsure{déterminé|déterminés|déterminées} \unsure{par} $\widehat L_F,\widehat L_{F'}$.}

% @ b), second condition
b) Pour que $X$ et $Y$ soient compatibles, il f. et s. que pour toute strate ouverte de $X$ et toute strate ouverte de $Y$, \deleted{on ait} celles-ci soient disjointes ou identiques.

% @ Cor., two lines
\uline{Cor.}
\[
\begin{aligned}
\text{\deleted{b)} a)}\quad &\widehat L_X\unsure{\ne|=}\widehat L_Y\ \deleted{\unsure{\Longrightarrow|\text{donc}}}\Longleftrightarrow\ X=Y\\
\text{b)}\quad &\widehat L_X\subset\widehat L_Y\ \deleted{\unsure{\widehat L_X|\widehat L_X\,\widehat L_X}}\ \text{\unsure{est} fermé dans }\widehat L_Y\ \Longleftrightarrow\ X\leq Y
\end{aligned}
\]

% @ Scholie, bottom
\uline{Scholie} On récupère $M$, avec $\ll$, $\unsure{R|\mathcal R}$ (compatibilité), en termes de $L$, et de l'ens \inserted{$\widehat M$} des $\widehat L_X$
\[
\widehat M=\operatorname{Im}\bigl(\underbrace{M\xrightarrow{X\mapsto\widehat L_X}\mathcal P(\mathcal P}_{\text{application \uline{injective}}}(L))\ \in\mathcal P(\mathcal P(\mathcal P(L)))
\]
